\documentclass[12pt,a4paper]{article}
\usepackage[utf8]{inputenc}
\usepackage[T1]{fontenc}
\usepackage[brazilian]{babel}
\usepackage{geometry}
\usepackage{hyperref}
\usepackage{listings}
\geometry{margin=2.5cm}
\title{Núcleo de Situação de Saúde (NSS)\\Sprint 1 -- Implementação}
\author{Zadoque Pires de Deus Souza Carneiro}
\date{Agosto de 2026}
\begin{document}
\maketitle

\section{Objetivo}
A Sprint 1 estabeleceu uma primeira versão executável do NSS. O ponto central foi integrar a fonte real do projeto, o SINAN, por meio da API atual do PySUS, e conectar a ingestão à arquitetura Medallion e à API FastAPI.

\section{Integração com o PySUS}
A implementação foi alinhada à referência oficial do PySUS. A API documenta a função de alto nível \texttt{pysus.api.sinan(disease, year, ...)} para obtenção dos registros SINAN. A função aceita códigos de doença como \texttt{DENG}, \texttt{TOXC} e \texttt{ZIKA}, além de um ano ou uma lista de anos.

O projeto usa explicitamente:
\begin{lstlisting}[language=Python]
from pysus.api import sinan

df = sinan(
    disease="DENG",
    year=2026,
    as_dataframe=True,
    show_progress=True,
)
\end{lstlisting}

O parâmetro \texttt{as\_dataframe=True} é importante porque a implementação atual do PySUS, por padrão, retorna caminhos dos arquivos baixados. Com esse parâmetro, o resultado é materializado como \texttt{pandas.DataFrame}. Assim, não foi utilizada a API antiga baseada em objetos \texttt{SINAN.get(...)}.

\section{Arquitetura Medallion}

\subsection{Bronze}
A Bronze recebe o DataFrame diretamente do PySUS e não aplica regras de negócio. O arquivo é salvo em Parquet com metadados de ingestão:
\begin{verbatim}
data/bronze/sinan/
  disease=DENG/source_year=2026/
  ingestion_date=YYYY-MM-DD/batch_id=YYYYMMDDTHHMMSSZ/
  data.parquet
  metadata.json
\end{verbatim}
O \texttt{metadata.json} registra doença, ano, lote, horário, quantidade de registros, colunas e fonte.

\subsection{Silver}
A transformação Silver seleciona as colunas necessárias ao primeiro caso de uso: \texttt{DT\_NOTIFIC}, \texttt{SEM\_NOT}, \texttt{NU\_ANO}, \texttt{ID\_MUNICIP}, \texttt{SG\_UF\_NOT}, \texttt{CLASSI\_FIN} e \texttt{EVOLUCAO}. A data é convertida para datetime, o código municipal é normalizado para sete dígitos, o código da UF para dois dígitos e são adicionadas sigla e nome da UF. Registros duplicados segundo as colunas do caso de uso são removidos.

\subsection{Gold}
A Gold deixa o dado pronto para consulta. Cada registro agregado contém:
\begin{itemize}
\item \texttt{disease\_code};
\item \texttt{year} e \texttt{month};
\item \texttt{cd\_uf} e \texttt{nm\_uf};
\item \texttt{cd\_mun} e \texttt{nm\_mun};
\item \texttt{cases\_total}.
\end{itemize}
O resultado é salvo em \texttt{data/gold/sinan/disease=<doenca>/year=<ano>/gold\_cases.parquet}.

O nome do município não é inventado a partir do código. Nesta sprint, \texttt{cd\_mun} é preservado e \texttt{nm\_mun} fica reservado para a futura dimensão territorial oficial do IBGE.

\section{Pipeline executável}
Foi criada a CLI \texttt{app.pipeline.run}, que executa sequencialmente:
\begin{verbatim}
PySUS -> Bronze -> Silver -> Gold
\end{verbatim}
Exemplo:
\begin{verbatim}
docker compose --profile pipeline run --rm pipeline \
  --disease DENG --year 2026
\end{verbatim}

\section{FastAPI}
A API foi ajustada para consultar somente a Gold materializada. Foram implementados:
\begin{itemize}
\item \texttt{GET /health};
\item \texttt{GET /diseases};
\item \texttt{GET /heatmap/uf?disease=DENG\&year=2026};
\item \texttt{GET /heatmap/municipios?disease=DENG\&cd\_uf=33\&year=2026}.
\end{itemize}
Também foi adicionada configuração CORS para permitir o consumo pelo front-end em \texttt{localhost:8080} e \texttt{127.0.0.1:8080}.

\section{Correção do erro 404}
O erro observado em \texttt{/heatmap/uf} não era uma rota inexistente. A implementação anterior levantava \texttt{404} quando não havia uma pasta/arquivo Gold para a doença solicitada. A causa era a ausência do dado processado no volume esperado.

A solução adotada foi separar claramente ingestão e consulta. A pipeline cria \texttt{data/gold}, e a API monta o mesmo diretório em \texttt{/data}. Quando a Gold ainda não existe, a API continua retornando \texttt{404}, mas agora informa o comando necessário para materializá-la. Isso evita mascarar ausência de dados com dados fictícios.

\section{Docker Compose}
O back-end possui dois serviços:
\begin{itemize}
\item \texttt{api}: FastAPI/Uvicorn;
\item \texttt{pipeline}: execução da pipeline PySUS sob demanda.
\end{itemize}
Ambos compartilham \texttt{../data:/data}. Assim, o Parquet criado pelo container da pipeline fica disponível para o container da API.

\section{Front-end}
O protótipo HTML consulta \texttt{/heatmap/uf} apenas após o usuário pressionar ``Consultar''. A consulta automática inicial foi removida para não gerar erros enquanto a Gold ainda não tiver sido criada.

\section{Critério de conclusão}
A Sprint 1 é considerada concluída quando uma doença/ano é processada com PySUS, os três níveis Bronze/Silver/Gold são materializados, a API consegue consultar a Gold via Docker Compose e o front-end consegue exibir os agregados por UF.

\section{Conclusão}
A principal mudança desta revisão foi substituir a suposição de uma API antiga do PySUS por sua interface de alto nível documentada atualmente. O projeto passa a possuir um fluxo real de dados, em vez de apenas um esqueleto de API esperando Parquets produzidos externamente.

\end{document}
